% ---------------------------------------------------------------------------
% DRAFT SEGMENT — the cue manipulation and how the donor is chosen
% Standalone; not \input anywhere yet. Stitch into Methods later.
% Code: hrtf/modify/donor_detail.py       (the composite)
%       hrtf/processing/midline.py        (applied to the measured arc)
%       hrtf/analysis/donor_selection.py  (vanopstal_dtfs / isim / composite_isim)
%
% Code and text now agree: shortlist()/select_donor() gate on ridge_slope <= 0.5
% and rank by |r_match - TARGET_R_MATCH|, both from pair_metrics() on the detail
% domain; DONOR_POOL holds the behaviourally qualified donors and
% MIN_DONOR_ELEVATION_GAIN is the criterion. isim()/vanopstal_dtfs() remain in
% the module as diagnostics and must NOT be reported as similarity.
% ---------------------------------------------------------------------------

\subsection{Spectral cue manipulation}

Participants were trained and tested with a modified version of their own
head-related transfer functions in which the coarse spectral envelope was their
own and the fine spectral detail was another listener's. For each direction and
ear, the log-magnitude spectrum was separated by cepstral truncation into an
envelope, retaining the first four coefficients, and a residual containing
everything finer. The modified spectrum was the participant's own envelope plus
a donor's residual. Own phase was retained and each direction was rescaled to
its original broadband energy, so interaural time differences and broadband
interaural level differences were unchanged and only spectral shape finer than
the envelope scale was altered.

Retaining four coefficients follows Kulkarni and Colburn
% \citet{kulkarni_role_1998}
, who smoothed transfer functions by cepstral truncation and found listeners
could no longer discriminate smoothed from original once roughly sixteen
coefficients were retained; four is well inside the range where the change is
audible and leaves the envelope too coarse to carry an elevation-dependent
pattern. Splitting the spectrum this way rather than substituting the donor's
transfer function whole was verified to decorrelate the spectral-to-spatial map
as thoroughly as a whole substitution, whereas the complementary
composite---donor envelope with own detail---left the native map intact. Detail
therefore carries the elevation cue and the envelope carries timbre, allowing
the cue to be replaced while level, head shadow and individual coloration are
retained.

The manipulation was applied to the measured median-plane transfer functions
and the full directional set was regenerated from them through the
spherical-head model used to build the unmodified set (Sec.~\ref{sec:hrir}), so
that native and modified sets share their interaural time differences exactly
and the envelope is fitted to the pinna response rather than to the pinna
response multiplied by the head shadow.

\subsection{Donor qualification}

Donor transfer functions were taken from a pool of previously recorded
participants. A recording qualified as a donor only if that participant had
themselves localized accurately with their own unmodified transfer function in
the same virtual environment, quantified as the elevation gain of their first
completed localization block (criterion: gain $\geq 0.8$; targets within
$\pm30^\circ$ elevation). The manipulation transplants one listener's spectral
detail to another, and this criterion establishes that the pattern being
transplanted is one a human auditory system demonstrably reads: it is evidence
that the cue works, obtained from the only instrument that can supply it.

The criterion is behavioral because the acoustic alternatives are not
predictive. Among the recordings for which both were available, the depth and
elevation dependence of a donor's own spectral detail---the natural acoustic
index of how much cue it carries---was unrelated to how well that participant
localized with it (Spearman $\rho = +0.18$, $n = 24$); the recording with the
most pronounced fine structure of the entire set belonged to a participant whose
elevation gain was $0.05$. Qualification is applied as an inclusion criterion
only. Accurate localization shows that a pattern is usable; inaccurate
localization is ambiguous between a weak cue, a rendering failure and an
inattentive listener, so an unqualified recording is treated as unproven rather
than as unsuitable.

\subsection{Donor selection}

Among qualified donors, the donor assigned to a participant was the one whose
composite came closest to an intended perturbation size while remaining
unabsorbable by that participant's existing spectral-to-spatial map. Both
quantities were computed between the participant's own and the modified spectral
detail on the median-plane arc, in the 5657--11314\,Hz band, averaged over ears.

Perturbation size was quantified as $r_\mathrm{match}$, the mean over elevations
of the correlation between the participant's own detail at that elevation and
the modified detail at the same elevation. This is the quantity that the
mechanism proposed by Van Wanrooij and Van Opstal
% \citet{vanwanrooij_relearning_2005}
---a correlation between the sensory input and a stored spectral representation
of the listener's own ears---predicts should govern performance. It was targeted
mid-range rather than minimized, because the perturbation is bounded from both
sides: in their monaural experiment, listeners whose cues were not sufficiently
perturbed never reached stable performance, while Trapeau et al.
% \citet{trapeau_fast_2016}
found that larger perturbations produced a greater initial loss and slower, less
complete recovery. The target was set to the median $r_\mathrm{match}$ between
qualified donors, so that the perturbation corresponds to a typical difference
between two people whose own cue is known to work.

Absorbability was quantified as the slope of the regression of the
best-matching own elevation on true elevation. A slope near unity means the
existing map can still read the modified set as a coherent elevation, displaced
rather than replaced, and can accommodate the manipulation as a constant
response bias; this is the configuration that Van Wanrooij and Van Opstal's
non-adapting listeners occupied. Candidates with a slope above $0.5$ were
excluded, and among the remainder the donor closest to the target
$r_\mathrm{match}$ was selected.

Both measures are required. Across all donor--listener pairings available here
they were nearly independent (Spearman $\rho = +0.14$), and of the pairings
whose $r_\mathrm{match}$ indicated a strong perturbation ($\leq 0.3$), $31\%$
nonetheless had a slope above the exclusion threshold---a large change in the
spectral cue that the existing map could still absorb.

Both were computed on the spectral detail rather than on whole transfer
functions. Correlating transfer functions requires the component common to all
directions to be removed by a reference derived from directions other than those
being compared, which for the spectral variability index
% \citep{trapeau_fast_2016}
is a diffuse-field average that the present recordings do not sample. A
reference formed from the median-plane arc itself cannot substitute, since it is
the mean of the same spectra that are then correlated: the residuals sum to zero
by construction, the mean pairwise correlation is fixed at $-1/(N-1)$, and the
resulting index is a constant (measured: $1.02 \pm 0.02$ across 31 recordings,
against $1.056$ predicted for $N = 19$ elevations). The cepstral residual is
subject to no such constraint, because the envelope removed from each direction
is a smoothed copy of that same direction and directions remain independent.

% --- open items ------------------------------------------------------------
% CITATIONS — verified 2026-08-18 unless marked:
%   Van Wanrooij & Van Opstal (2005) J Neurosci 25(22):5413-5424,
%     doi:10.1523/JNEUROSCI.0850-05.2005. "Relearning sound localization with a
%     new ear" (SINGULAR -- the MONAURAL perturbation study). Do NOT confuse
%     with Hofman, Van Riswick & Van Opstal (1998) Nat Neurosci, "with new ears"
%     (plural), the BINAURAL mold study. I_sim definition = their Materials and
%     Methods, "DTF similarity index"; DTF definition = their Eq. 1; the 0.5 and
%     0.3 anchors = their Fig. 7A and 7C; r^2 = 0.94 = their Fig. 6 (n=6).
%     All quoted from the full text, verified.
%   Kulkarni & Colburn (1998) Nature -- cepstral smoothing, discrimination
%     threshold near 16 coefficients. [volume/pages NOT VERIFIED]
%   Trapeau et al. (2016) JASA -- VSI. Cited only if the paragraph declining to
%     use it names it. [volume/pages NOT VERIFIED]
%
% - QUALIFIED POOL = 7 (CO, AGV, SW, VD, AH, FS, FD), from 27 recordings with an
%   own-HRTF block. Every one of the 12 listeners resolves to an admissible
%   donor: 6 in band, 6 widened, no fallbacks. Target r_match 0.58 = median over
%   the 21 qualified-donor pairs (min 0.32, Q1 0.51, Q3 0.66, max 0.77).
%   RECOMPUTE THE TARGET WHENEVER THE POOL CHANGES -- it is defined from the pool.
% - 0.8 falls in the gap between FD at 0.82 and AS at 0.63, so it is not making a
%   fine distinction. State it as a cut in a gap, not as a graded criterion.
% - GS IS A PROBLEM CASE: the only admissible qualified donor for GS is AGV at
%   r_match 0.90, i.e. almost no perturbation. Either GS needs a donor that is
%   not yet qualified, or GS cannot be run under this rule. Check before
%   scheduling GS.
% - Five candidates still have no own-HRTF block (GLK, IM, MSc, NW, UG) -- they
%   have other runs, just none with their unmodified HRTF. One ~60-trial block
%   each would qualify them and widen the pool.
% - CORRECTED: an earlier version of this file called the EG distribution
%   bimodal with a suspicious cluster at zero. That came from
%   vsi_rmse.find_participants() missing 6 subjects (it searched only the flat
%   results/pilot/<id>.pkl and not the nested results/pilot/<id>/<id>.pkl, and
%   preferred the pilot pair where an id existed in both trees, which shadowed
%   the active AS and PF). Both are fixed; with the fuller set the spread is
%   closer to continuous. Do not repeat the bimodality claim.
% - WHY NO I_SIM. Van Opstal's index was tested and rejected for this role: it is
%   the SD down each column of the correlation matrix, and SD is invariant to
%   permutation of the rows, so it cannot distinguish a listener's own ear from
%   their own ear relabelled by 17 deg (measured: identical value 0.542 for
%   shifts of 0, 4, 8 and 17 deg, while the correlation at matching elevation
%   fell 1.00 -> 0.08). It also scores a cue-REMOVED ear (own envelope only,
%   0.395) as less perturbed than a real donor composite (0.272), so it does not
%   separate replacement from removal. Keep this in the file; it is the answer if
%   a reviewer asks why the obvious measure was not used.
% - Detail strength has no repeatability estimate (every arc is measured once),
%   so do not report it to more precision than the ranking requires, and do not
%   claim a difference between adjacent donors.
% - NO REPEATABILITY ESTIMATE exists for any of these measures: every arc is
%   measured once. Re-recording one subject would fix this cheaply, and two
%   recordings HAVE been remade (AS 2026-08-18, GS 2026-07-28) -- but as
%   replacements, not repeats, so they do not serve.
% - Sec.~\ref{sec:hrir} is a placeholder label -- point it at whichever section
%   describes the HRIR recording and the spherical-head azimuth expansion.
% - Regenerate if anything changes: the own-ear and between-listener I_sim
%   values and their n, and the target once set.
